% !TEX root = ../notes_sovereignai.tex
% Chapter 5: Second Brain. Obsidian on the local engine.
%%%%%%%%%%%%%%%%%%%%%%%%%%%%%%%%%%%%%%%%%%%%%%%%%%%%%%%%%%%%%%%%%%%%%%

\index{Obsidian}\index{embeddings}\index{retrieval}\index{second brain}\index{markdown}
\chapter{Second Brain}
\printchapterversion{05_secondbrain}
\newthought{Chat With What You Know.}
\vspace{1.5em}

% ==============================================================================================
% THE WHY
% ==============================================================================================

\section{The Why}\index{retrieval augmented generation}\index{RAG|see{retrieval augmented generation}}

\newthought{The model answers from what it was trained on}, not from what you have lived. It has
never read your meeting notes, your reading highlights, the half-finished thoughts in
your vault. Point your notes at the engine and it answers from your own record
rather than from general knowledge, while the notes themselves never leave your
machine.

\marginnote{\newthought{Obsidian}\index{Obsidian!vault} keeps a vault as plain markdown files in a folder
on your own disk, with no server and no account behind it, which is what makes a
private path from the vault to the engine a matter of pointing a plugin at an
address.}

\newthought{The tool is Obsidian} with two community plugins\index{Obsidian!plugins}: one to chat with the
vault, one to find related notes by meaning. Both point at the Chapter~2 engine over
the private network from Chapter~3, so a question typed on your laptop is answered by
the rig at home, grounded in your writing.

\newthought{By the end of this chapter} you will have:
\begin{itemize}
    \item pulled an embedding model and exposed it through the engine;
    \item connected Obsidian's chat and semantic-search plugins to the rig;
    \item asked a question and received an answer drawn from your own notes.
\end{itemize}

% ==============================================================================================
% NOTES AS CONTEXT
% ==============================================================================================

\section{Notes as Context}\index{embedding model}\index{semantic search}\index{context window}

\newthought{A model reads only its prompt.} To answer from your notes, the relevant
notes must be found and placed into the prompt before the model runs. That is
retrieval-augmented generation\index{retrieval augmented generation}: search the vault for what bears on the question, hand
those passages to the model as context, and let it answer over them. The model is not
retrained; it is briefed.

\marginnote{\newthought{Embeddings}\index{vector} turn each note into a vector, a point in meaning
space, so that notes about the same idea land near each other even when they share no
words. A question becomes a vector too, and the nearest notes are its answer's
sources. This is what makes search find sense, not just strings.}

\newthought{Two models, two jobs.} Search needs an embedding model, small and fast,
that vectorises every note once and the question each time; answering needs the chat
model from Chapter~2. Both run on the rig through Ollama, so the embedding model is a
one-line pull and the vault is indexed against your own hardware.

\marginnote{\newthought{Already partly solved.}\index{knowledge base} Open WebUI carries its own retrieval:
drop a PDF or a folder into its knowledge feature and it embeds and searches them the
same way. Obsidian is the version that lives inside your notes;
the next chapter points at the broader cases.}

\newthought{Where the work runs.}\index{tailnet} Obsidian runs on your laptop, the engine on the
rig. The plugins reach the engine at its address on your tailnet, so indexing and
chat happen on the 3090 while the vault stays on the laptop, and nothing crosses the
public internet.

% ==============================================================================================
% WIRE OBSIDIAN
% ==============================================================================================

\section{Wire Obsidian}\index{Copilot}\index{Smart Connections}

\newthought{Pull the embedding model}\index{nomic-embed-text} on the rig, alongside the chat model:

\begin{verbatim}
$ ollama pull nomic-embed-text        # small, fast, made for embeddings
$ ollama list                         # qwen3:14b and nomic-embed-text present
\end{verbatim}

\marginnote{\newthought{Reach over the tailnet.} Use the rig's tailnet address, not
localhost, since Obsidian runs on the laptop. The mesh from Chapter~3 is the
authentication, so the plugins can speak to Ollama directly and safely.}

\newthought{Install the plugins.} In Obsidian, enable community plugins and add
Copilot, for chat with the vault, and Smart Connections, for semantic search. In each
plugin's settings, choose Ollama as the provider and set the endpoint to the rig:

\begin{verbatim}
Provider ......... Ollama (OpenAI-compatible)
Base URL ......... http://rig.your-tailnet.ts.net:11434
Chat model ....... qwen3:14b
Embedding model .. nomic-embed-text
\end{verbatim}

\marginnote{\newthought{If the plugin cannot connect}\index{cross-origin requests}, set \texttt{OLLAMA\_ORIGINS}
on the rig so the engine accepts requests from the Obsidian app, the one piece of
cross-origin housekeeping these plugins need.}

\newthought{Index the vault}\index{vector}\index{embeddings}, a one-time pass that embeds every note; on a normal
vault it finishes in minutes and updates incrementally afterwards. Then open the chat
panel and ask something only your notes know.

\newthought{What you now own} answers from your own record. You ask ``what did I
conclude about the rig's power budget?'' and the assistant answers in your own
words, citing the note you wrote while building it, with the reasoning done on the
3090 and the vault never leaving your laptop.

% ==============================================================================================
% BRIDGE
% ==============================================================================================

\newthought{The engine now answers from your own record}, not only from its
training, and chat, terminal, voice, and notes all run on the one box you own. The
stack is complete enough to live on. The last chapter is a map rather than a build:
where to take the rig when you want more from it.

\marginnote{\newthought{feedback}}
